\documentclass[xcolor=svgnames,11pt]{beamer}
% \documentclass[handout,xcolor=svgnames,11pt]{beamer}

\usepackage{graphicx}

\usepackage[italian, english]{babel}
\usepackage[utf8x]{inputenc}
\usepackage[T1]{fontenc}
\usepackage[scale=0.9,sfdefault,light]{roboto}
\usepackage{hyperref}
\usepackage{tabularx}
\usepackage{url}
\usepackage{pbox}

\usepackage{etoolbox}
\makeatletter
\patchcmd{\insertverticalnavigation}%
{\ifx\beamer@nav@css\beamer@hidetext{\usebeamertemplate{section in sidebar}}\else{\usebeamertemplate{section in sidebar shaded}}\fi}%
{{\usebeamertemplate{section in sidebar}}}{}{}
\makeatother

\urlstyle{sf}

\usetheme{Goettingen}
\definecolor{leorange}{HTML}{FBA100}
\definecolor{leblue}{HTML}{003A70}
\definecolor{myblack}{RGB}{14,14,14}
\definecolor{lelightblue}{HTML}{BFCDDB}

\usecolortheme[named=leblue]{structure}
\beamertemplatenavigationsymbolsempty

\beamertemplateshadingbackground{white!15}{leorange!5}
\setbeamertemplate{sidebar canvas right}[vertical shading]%
[top=lelightblue,bottom=lelightblue]

\mode<presentation>{\usebackgroundtemplate{\includegraphics
    [height=\paperheight]{img/back_le.png}}}

\setbeamertemplate{blocks}[rounded][shadow=false]
\setbeamercolor{block body}{bg=leorange!50}
\setbeamercolor{block title}{bg=leorange!80}

\setbeamercovered{transparent}

\title{\textbf{Let's encrypt! Certificati SSL per tutti!}}
\author{Nicola Corti}
\institute{Gruppo Utenti Linux Pisa \\ \medskip \includegraphics[height=1.5cm]{img/gulp_logo.png}}
\logo{img/gulp_logo.png}
\date{03 Febbraio 2016}


\begin{document}

\begin{frame}
	\titlepage
\end{frame}

\section{Intro}
\begin{frame}{}
\begin{center}
\begin{Huge}
\textcolor{leorange}{\emph{Introduction}}
\end{Huge}
\end{center}
\end{frame}



\begin{frame}{Cosa è Let's encrypt?}
\begin{center}
\includegraphics[width=5cm]{img/logo_le.png}
\end{center}

\begin{block}{}
Let's encrypt è una \textbf{Certification Authority (CA)} che rilascia
certificati SSL/TLS a titolo \textbf{gratuito}.
\end{block}
\pause \medskip
\begin{itemize}
  \item Dal 5 Dicembre 2015 si trova in \textbf{Public Beta}.
  \pause
  \item Ad oggi ha rilasciato più di \textbf{480 k} certificati.
  \pause
  \item Fa dell'\textbf{automazione} il suo punto di forza.
  \pause
  \item \textbf{\url{https://letsencrypt.org/}}
\end{itemize}
\end{frame}

\subsection{Sponsors}
\begin{frame}{Sponsors}
\begin{center}
Chi sta supportando il progetto?
\medskip
\vspace{0.5cm}
\includegraphics[width=10cm]{img/sponsors.png}
\end{center}
\end{frame}

\subsection{History}
\begin{frame}{History}
\begin{description}
  \item[2012] Iniziato il progetto dentro \textbf{Mozilla}
  \pause
  \item[11-2014] Let's Encrypt \textbf{annunciato pubblicamente}
  \pause
  \item[01-2015] Protocollo \textbf{ACME} sottoposto all'\textbf{IETF} per essere standardizzato
  \pause
  \item[04-2015] \textbf{ISRG} (Internet Security Research Group) e la \textbf{Linux Foundation}
  si uniscono al progetto
  \pause
  \item[09-2015] Rilasciato il primo certificato per \textbf{\href{https://helloworld.letsencrypt.org/}{helloworld.letsencrypt.org}}
  \pause
  \item[10-2015] Il certificato intermedio di L.E. diventa \textbf{\emph{cross-signed}} grazie a IdenTrust.
  (\emph{Let's Encrypt is trusted!})
  \pause
  \item[12-2015] \textbf{Public beta!}
\end{description}
\end{frame}

\subsection{Why?}
\begin{frame}{Why?}
\begin{itemize}
  \item \textbf{Free}
  \medskip\pause
  \item \textbf{Automatico}
  \medskip\pause
  \item \textbf{Open}
  \medskip\pause
  \item \textbf{Sicuro}
  \medskip\pause
  \item \textbf{Trasparente}
\end{itemize}
\end{frame}

\section{Server}
\begin{frame}{}
\begin{center}
\begin{Huge}
\textcolor{leorange}{\emph{Server}}
\end{Huge}
\end{center}
\end{frame}



\begin{frame}{Boulder}
\begin{block}{Architettura}
Il sistema di Let's Encrypt si basa su \textbf{3 Componenti}: un \textbf{server},
un \textbf{client} e il \textbf{protocollo} che permette a server e client di comunicare.
\end{block}

\medskip \pause

Il server si chiama \textbf{Boulder} ed è scritto interamente in \textbf{Go}.
Si occupa di gestire tutte le procedure per il \textbf{rilascio} e \textbf{revoca} dei certificati.

\medskip \pause

Di fatto si tratta di un server HTTP che espone un'interfaccia \textbf{RESTful}

\medskip \pause

\begin{center}
\visible<4>{\includegraphics[width=0.3cm]{img/logo_gh.pdf}} \textbf{\url{https://github.com/letsencrypt/boulder}}
\end{center}
\end{frame}

\section{Client}
\begin{frame}{}
\begin{center}
\begin{Huge}
\textcolor{leorange}{\emph{Client}}
\end{Huge}
\end{center}
\end{frame}



\begin{frame}{letsencrypt}
Il client si chiama (semplicemente) \textbf{letsencrypt} ed è scritto interamente in \textbf{Python}.
Si occupa di farvi interagire con il server remoto e \textbf{gestisce i vostri certificati}.
\pause
\begin{itemize}
  \item Installabile tramite il pacchetto \textbf{letsencrypt} (solo debian \emph{sid/stretch}).
  \item Installabile tramite \textbf{clone} del repository.
\end{itemize}
\pause
\begin{center}
\visible<3>{\includegraphics[width=0.3cm]{img/logo_gh.pdf}} \textbf{\url{https://github.com/letsencrypt/letsencrypt}}
\end{center}
\end{frame}

\begin{frame}{Plugins}
Lo scopo principale del client è quello di \textbf{semplificare e automatizzare} tutto il processo di
autenticazione e di creazione del certificato.

\medskip\pause

Per questo motivo sono stati sviluppati dei plugin per il setup automatico dei certificati ottenuti
sui principali web browser: \textbf{apache} ed \textbf{nginx}.

\visible<2>{
\begin{center}
\includegraphics[height=2cm]{img/apache_logo.png}
\includegraphics[height=2cm]{img/nginx_logo.png}
\end{center}}

\end{frame}

\section{Protocol}
\begin{frame}{}
\begin{center}
\begin{Huge}
\textcolor{leorange}{\emph{Protocol}}
\end{Huge}
\end{center}
\end{frame}


\begin{frame}{ACME}

\begin{block}{}
Il protocollo che utilizzano client e server per interagire si chiama \textbf{Automated Certificate Management Environment (ACME)}.
\end{block}

\medskip \pause
ACME si base sullo scambio di file \textbf{JSON} firmati (anche detti \textbf{JWS, Json Web Signature}).
Questi documenti contengono le richieste inviate dal client e le risposte ottenute dal server.

\medskip \pause
Lo scambio di questi documenti deve avvenire necessariamente attraverso \textbf{HTTPS}.
\end{frame}

\begin{frame}{ACME}
\begin{center}
\includegraphics[width=9.5cm]{img/acme_spec.png}

\includegraphics[width=0.3cm]{img/logo_gh.pdf} \textbf{ \url{https://github.com/letsencrypt/acme-spec}}
\end{center}
\end{frame}


\begin{frame}{ACME}
Il protocollo ACME ci permette di eseguire tre operazioni:
\pause\medskip
\begin{enumerate}
  \item \textbf{Provare} che noi siamo i responsabili del dominio \emph{example.com}
  \pause\medskip
  \item \textbf{Ottenere} un certificato per il dominio \emph{example.com}
  \pause\medskip
  \item \textbf{Revocare} o \textbf{Rinnovare} un certificato per il dominio \emph{example.com}
\end{enumerate}
\end{frame}

\subsection{Security}
\begin{frame}{Security}
Le interazioni fra client e server sono cifrate tramite una coppia di \textbf{chiavi asimmetriche}
(una \textbf{privata} ed una \textbf{pubblica}) che viene generata durante il primo avvio del client.

\medskip \pause

\begin{block}{}
Il server ci permette di \textbf{provare} che siamo gli amministratori del dominio desirato, inviandoci
una o più \textbf{challenge} che dobbiamo risolvere.
\end{block}

\medskip \pause

Ogni interazione con il server è inoltre corredata da un \textbf{nonce} che permette di evitare attacchi
di tipo \textbf{Replay}.

\end{frame}

\subsection{Challenges}
\begin{frame}[fragile]{Challenges}
\begin{center}

Il server può, a sua discrezione, inviare una o più challenges.
\medskip\pause

\begin{tabular}{p{3cm}p{6cm}}

\hline
\textbf{Tipo} & \textbf{Descrizione} \\
\hline
\textcolor{leblue}{\textbf{\small Simple HTTP}} & \textcolor{leblue}{{\small Si chiede di posizionare un \textbf{file} all'interno del
 proprio webserver contenente un \emph{token} specifico (Accettato sia HTTP che HTTPS).}} \\
\hline
\textbf{\small DNS} & {\small Si chiede di inserire un \textbf{record TXT} all'interno del proprio DNS.} \\
\hline\hline
\textbf{\scriptsize Proof of possession} & {\scriptsize Si chiede di \textbf{firmare un documento} utilizzando una chiave
che il server può ricondurre al client.} \\
\hline
\textbf{\scriptsize Domain Validation with Server Name Indication} & {\scriptsize Si chiede di \textbf{configurare un server
TLS} ad uno specifico indirizzo IP (tramite un record A nel DNS).} \\
\hline
\end{tabular}
\end{center}
\end{frame}

\subsection{The protocol}
\begin{frame}{Domain Validation}
\begin{center}
\includegraphics[width=9cm]{img/proto_1.png}
\end{center}
\end{frame}

\begin{frame}{Domain Validation}
\begin{center}
\includegraphics[width=9cm]{img/proto_2.png}
\end{center}
\end{frame}

\begin{frame}{Certificate Issuance}
\begin{center}
\includegraphics[width=9cm]{img/proto_3.png}
\end{center}
\end{frame}

\subsection{Cert revocation}
\begin{frame}{Certificate Revocation}
\begin{center}
\includegraphics[width=9cm]{img/proto_4.png}
\end{center}
\end{frame}

\section{Certs}
\begin{frame}{}
\begin{center}
\begin{Huge}
\textcolor{leorange}{\emph{Certificates}}
\end{Huge}
\end{center}
\end{frame}


\subsection{DV}
\begin{frame}{DV Certificates}
\begin{block}{DV}
I certificati rilasciati sono di tipo \textbf{Domain Validated}. Provano cioè
che si è effettivamente i proprietari di un determinato dominio.
\end{block}
\medskip\pause

Certificati di tipo \textbf{Organization Validation} oppure \textbf{Extended Validation}
richiederebbero verifiche esplicite del soggetto che fa richiesta del certificato.

\visible<2>{
\begin{center}
\includegraphics[width=6cm]{img/dv_example.png}

\includegraphics[width=6cm]{img/ev_example.png}
\end{center}}
\end{frame}

\subsection{Cross Signing}
\begin{frame}{Cross Signing}
Tutti i certificati rilasciati sono \emph{Cross-signed} da parte di \textbf{IdenTrust}.
Questo permette ai certificati di L.E. di essere riconosciuti dalla maggioranza dei browser.

\medskip\pause

In questo modo possiamo evitare errori quali:

\visible<2>{
\begin{center}
\includegraphics[width=7cm]{img/cert_error.png}
\end{center}}
\end{frame}

\begin{frame}{Cross Signing}
\begin{center}
\includegraphics[width=10cm]{img/certs.png}
\end{center}
\end{frame}

\begin{frame}{Validity}
I certificati hanno una validità di \textbf{90 giorni} dopo i quali scadono, e non sono
considerati più validi dai browser.

\medskip\pause

\begin{block}{\emph{90-days validity}}
Questa non è una novità nel mondo della sicurezza. \textbf{Rinnovare i certificati frequentemente}
è una buona abitudine e permette di \textbf{limitare i danni} di eventuali certificati compromessi.
\end{block}

\medskip\pause
Riceveremo delle mail di \emph{remind} che ci ricordano della scadenza.
È possibile automatizzare il processo di rinnovo tramite dei task di \textbf{cron}.
\end{frame}


\section{Setup}
\begin{frame}{}
\begin{center}
\begin{Huge}
\textcolor{leorange}{\emph{Setup}}
\end{Huge}
\end{center}
\end{frame}


\subsection{Requirements}
\begin{frame}{Requirements}
Per eseguire il client è necessario:
\begin{itemize}
  \item Essere su un sistema \textbf{Unix like}.
  \item Avere \textbf{Python 2.6} o \textbf{2.7}.
  \item Avere accesso di \textbf{root} sul proprio sistema.
\end{itemize}

\medskip\pause

Inoltre, il plugin di autoconfigurazione di \textbf{apache}, funziona solamente con
OS Debian based: \textbf{Ubuntu 12.04+} e \textbf{Debian 7+}
\end{frame}

\subsection{Download}
\begin{frame}[fragile]{Download}
Per prima cosa è necessario scaricare il client \textbf{letsencrypt}.

\medskip\pause
Per Debian \emph{sid/stretch}:
\begin{block}{}
\begin{scriptsize}
\begin{verbatim}
$ sudo apt-get install letsencrypt
$ letsencrypt --help
\end{verbatim}
\end{scriptsize}
\end{block}

\medskip\pause
Per tutti gli altri:
\begin{block}{}
\begin{scriptsize}
\begin{verbatim}
$ git clone https://github.com/letsencrypt/letsencrypt
$ cd letsencrypt
$ ./letsencrypt-auto --help
\end{verbatim}
\end{scriptsize}
\end{block}

\medskip\pause

Da ora in poi ci riferiremo ai comandi con \textbf{letsencrypt-auto} assumendo che
si scelga di clonare il repository.

\end{frame}

\subsection{Run}
\begin{frame}[fragile]{Run}
Eseguire il client è molto semplice. È sufficiente invocare il comando:
\begin{block}{}
\begin{scriptsize}
\begin{verbatim}
$ ./letsencrypt-auto
\end{verbatim}
\end{scriptsize}
\end{block}
E verremo guidati durante il processo di issuing dei certificati.
\end{frame}

\begin{frame}[fragile]{Apache}
Se si vuole configurare \textbf{Apache} per utilizzare i certificati che vengono
generati, è possibile utilizzare il comando
\begin{block}{}
\begin{scriptsize}
\begin{verbatim}
$ ./letsencrypt-auto --apache -d example.com -d www.example.com
\end{verbatim}
\end{scriptsize}
\end{block}
Indicando con \textbf{-{}-apache} che si vuole configurare Apache e con \textbf{-d}
i vari domini per cui si sta richiedendo un certificato.
\end{frame}

\begin{frame}[fragile]{Contacts}
Durante la prima interazione ci verrà chiesta la nostra \textbf{email} per i contatti
e ci verrà chiesto di accettare i \textbf{Terms of service}.

\medskip\pause

È possibile saltare questi passaggi tramite i seguenti flag

\medskip

\begin{block}{}
\begin{scriptsize}
\begin{verbatim}
$ ./letsencrypt-auto --email admin@example.com --agree-tos
\end{verbatim}
\end{scriptsize}
\end{block}
\end{frame}

\subsection{Plugins}
\begin{frame}{Plugins}
\begin{center}
\begin{tabular}{p{2cm}p{0.5cm}p{0.5cm}p{5cm}}
\hline
\textbf{Plugin} & \textbf{A} & \textbf{I} & \textbf{Descrizione}\\
\hline
Apache & Y & Y & {\scriptsize Richiede ed installa automaticamente i certificati su Apache 2.4 (Debian based)} \\
Standalone & Y & N & {\scriptsize Avvia un webserver \emph{standalone} sulla porta 80/443} \\
Webroot & Y & N & {\scriptsize Ottiene il certificato andando a scrivere all'interno della \emph{root} di web server già attivo}\\
Manual & Y & N & {\scriptsize Fornisce i comandi da eseguire per ottenere un certificato su un'altro web server} \\
Nginx & Y & Y & {\scriptsize Richiede ed installa automaticamente i certificati su Nginx (sperimentale)} \\
\hline
\end{tabular}
\end{center}
\end{frame}


\begin{frame}[fragile]{certonly}
È possibile utilizzare i plugin con il supporto per l'\emph{authentication} (Y nella colonna A)
per ottenere \textbf{solamente i certificati} senza effettuare l'installazione degli stessi.

\medskip\pause

È sufficiente indicare l'opzione \textbf{certonly} nella linea di comando.

\medskip

\begin{block}{Standalone example}

Ad esempio tramite il comando
\begin{scriptsize}
\begin{verbatim}
$ ./letsencrypt-auto --standalone-supported-challenges \
 http-01 certonly -d example.com
\end{verbatim}
\end{scriptsize}

Si avvia un server standalone sulla porta 80 e si otterrà il certificato per il dominio example.com
\end{block}
\end{frame}

\begin{frame}{/etc/letsencrypt}
Il client andrà a salvare i certificati e le chiavi dentro la cartella \textbf{/etc/letsencrypt}.

\medskip\pause
{
\setbeamercolor{block body}{bg=red!50}
\begin{block}{}
Dentro questa cartella andranno a finire sia i \textbf{certificati} che le \textbf{chiavi pubbliche/private}.
È quindi \textbf{FONDAMENTALE} fare un backup della cartella e mantenerla in un \textbf{luogo sicuro}.
\end{block}
}
\medskip\pause
I certificati si troveranno dentro la cartella \textbf{/etc/letsencrypt/live/example.com/}, sono link
simbolici che verranno aggiornati ad ogni rinnovo dei certificati.
\end{frame}

\begin{frame}{/etc/letsencrypt}
I file che andremo a trovare sono:
\begin{description}
  \item[\textbf{\small privkey.pem}] Chiave privata del certificato. Da non condividere MAI!
  \item[\textbf{\small cert.pem}] Certificato del server (quello che viene inviato al client).
  \item[\textbf{\small chain.pem}] Elenco dei certificati intermedi delle C.A. collegati al certificato.
  \item[\textbf{\small fullchain.pem}] Elenco dei certificati intermedi delle C.A. collegati al certificato, più il certificato stesso.
\end{description}
\end{frame}

\subsection{Revoke}
\begin{frame}[fragile]{Revoke}
Per revocare un certificato è sufficiente invocare l'opzione \textbf{revoke}.

\medskip\pause

\begin{block}{}
\begin{scriptsize}
\begin{verbatim}
$ /.letsencrypt-auto revoke --cert-path example-cert.pem
\end{verbatim}
\end{scriptsize}
\end{block}
\end{frame}

\subsection{Renewal}
\begin{frame}[fragile]{Renewal}
L'operazione di rinnovo è molto semplice, e può essere effettuata semplicemente
invocando il comando \textbf{letsencrypt}, senza parametri.

\medskip\pause

È inoltre possibile impostare flag quali \textbf{-{}-renew-by-default} per effettuare
il rinnovo senza nessuna interazione da parte dell'utente.

\medskip\pause

In questo modo è possibile schedulare un task su \textbf{cron} per il rinnovo
automatico dei certificati prima dei 90 giorni.

\end{frame}

\subsection{Update}
\begin{frame}[fragile]{Update}
È importante tenere il client aggiornato, in quanto si tratta di una beta pubblica.


\medskip\pause
Per Debian \emph{sid/stretch}:
\begin{block}{}
\begin{scriptsize}
\begin{verbatim}
$ apt-get update && apt-get upgrade
\end{verbatim}
\end{scriptsize}
\end{block}

\medskip\pause
Per tutti gli altri:
\begin{block}{}
\begin{scriptsize}
\begin{verbatim}
$ cd letsencrypt
$ git pull
\end{verbatim}
\end{scriptsize}
\end{block}

\end{frame}

\begin{frame}[fragile]{General usage}
Nella stragrande maggioranza dei casi, sarà sufficiente invocare il comando
\textbf{letsencrypt} e farsi guidare dall'interfaccia grafica da terminale.

\medskip\pause

Let's Encrypt ci seguirà durante tutti i processi di generazione/rinnovo
di ogni certificato. Sarà sufficiente rispondere alle domande che ci vengono
chieste.

\medskip\pause

\textbf{Troppi parametri} da ricordare renderebbero poco snello il processo automatico
di letsencrypt, ma al tempo stesso permettono molta \textbf{flessibilità} e permettono
di includere letsencrypt dentro \textbf{script}.

\end{frame}

\subsection{Screens}
\begin{frame}{Screens}
\begin{center}
\includegraphics[width=9cm]{img/screen1.png}
\end{center}
\end{frame}
\begin{frame}{Screens}
\begin{center}
\includegraphics[width=9cm]{img/screen2.png}
\end{center}
\end{frame}
\begin{frame}{Screens}
\begin{center}
\includegraphics[width=9cm]{img/screen3.png}
\end{center}
\end{frame}
\begin{frame}{Screens}
\begin{center}
\includegraphics[width=9cm]{img/screen4.png}
\end{center}
\end{frame}
\begin{frame}{Screens}
\begin{center}
\includegraphics[width=9cm]{img/screen5.png}
\end{center}
\end{frame}
\begin{frame}{Screens}
\begin{center}
\includegraphics[width=9cm]{img/screen6.png}
\end{center}
\end{frame}
\begin{frame}{Screens}
\begin{center}
\includegraphics[width=9cm]{img/screen7.png}
\end{center}
\end{frame}
\begin{frame}{Screens}
\begin{center}
\includegraphics[width=9cm]{img/screen8.png}
\end{center}
\end{frame}


\section{Conclusions}
\begin{frame}{}
\begin{center}
\begin{Huge}
\textcolor{leorange}{\emph{Conclusions}}
\end{Huge}
\end{center}
\end{frame}



\subsection{Drawbacks}
\begin{frame}{Drawbacks}
\begin{itemize}
  \item Nessun supporto per \textbf{Organization Validation} o \textbf{Extended
  Validation}, troppo complessi da automatizzare.
  \medskip\pause
  \item Nessun supporto per le \textbf{wildcards} (*.example.com), forse in futuro.
  \medskip\pause
  \item Challenges solo tramite HTTP (public beta), supporto DNS ancora assente.
\end{itemize}
\end{frame}

\subsection{Resources}
\begin{frame}{Resources}
\begin{small}
\begin{itemize}
  \item How it works \textbf{\url{https://letsencrypt.org/howitworks/}}
  \medskip
  \item Tech details \textbf{\url{https://letsencrypt.org/howitworks/technology/}}
  \medskip
  \item Read the docs \textbf{\url{https://letsencrypt.readthedocs.org/}}
  \medskip
  \item Community board \textbf{\url{https://community.letsencrypt.org/}}
  \medskip
  \item Code \textbf{\url{https://github.com/letsencrypt/}}
  \medskip
  \item Mailing lists
  \begin{itemize}
    \item Client \textbf{\url{https://groups.google.com/a/letsencrypt.org/forum/\#!forum/client-dev}}
    \item Server \textbf{\url{https://groups.google.com/a/letsencrypt.org/forum/\#!forum/ca-dev}}
    \item ACME (IETF) \textbf{\url{https://www.ietf.org/mailman/listinfo/acme}}
  \end{itemize}
\end{itemize}
\end{small}
\end{frame}

\begin{frame}{}
\begin{center}
\begin{Huge}
{\color{leorange} \textbf{Domande...?}}
\end{Huge}

\vspace{1.5cm}
% Slides realizzate da:\\
% \textbf{Nicola Corti - corti.nico [at] gmail [dot] com}\\

\begin{center}
\begin{tabular}{>{\centering\arraybackslash}p{4cm}}
\includegraphics[height=0.2cm]{img/logo_web.pdf} \textbf{\href{https://ncorti.com}{ ncorti.com}} \\
\includegraphics[height=0.2cm]{img/logo_twitter.pdf} \textbf{\href{https://twitter.com/cortinico}{ @cortinico}} \\
\includegraphics[height=0.2cm]{img/logo_gh.pdf} \textbf{\href{https://github.com/cortinico}{ @cortinico}} \\
\end{tabular}
\end{center}

\bigskip

\begin{footnotesize}
Slides realizzate con \LaTeX\ Beamer.\\
La seguente presentazione \`e rilasciata sotto licenza\\
\textbf{Creative Commons - Attributions, Non Commercial, Share-alike}.
\end{footnotesize}
\\
\medskip
\includegraphics[height=0.5cm]{img/cc.png}

\end{center}
\end{frame}

\end{document}
